\section{Background: autonomy, coding agents and Claude Code}
\label{sec:background}

\subsection{Levels of autonomy and degrees of agency}

The idea that automation comes in degrees predates LLMs. \citet{parasuraman2000model} model
automation along four information-processing stages---information acquisition, analysis, decision
selection and action implementation---each of which can be automated to a different level. A
consultative LLM automates acquisition and analysis and leaves decision and action to the human. A
coding agent automates all four.

Recent frameworks transpose this view to LLM agents. \citet{feng2025levels} define five levels of
autonomy by the role the user takes---operator, collaborator, consultant, approver and observer---and
argue that autonomy is a \emph{design decision}, separable from capability. The distinction matters
here: the same model can be deployed as an advisor or as an operator, and goal drift is a property of
the deployment as much as of the model. \citet{kasirzadeh2025characterizing} characterise agents
along autonomy, efficacy, goal complexity and generality. Coding agents score high on all four:
they act directly on a codebase, pursue compound goals and are general-purpose.
\citet{chan2023harms} identify four characteristics that make an algorithmic system more agentic:
underspecification, directness of impact, goal-directedness and long-term planning. Each one
plausibly favours sub-task accumulation. Underspecified goals invite the agent to fill gaps with its
own sub-goals; direct impact means those sub-goals change the repository; long-term planning means
they persist in the context.

Two further positions frame the normative stakes. \citet{mitchell2025fully} argue that risks to people
grow with the control a user cedes, and that fully autonomous agents should not be developed.
\citet{sapkota2025agents} distinguish single-task, tool-using \emph{AI agents} from \emph{agentic AI}
systems characterised by dynamic task decomposition, multi-agent collaboration and persistent memory.
Claude Code, with its sub-agents, task lists and memory files, sits between the two. Finally,
\citet{chan2024visibility} argue that once agents act, \emph{visibility}---agent identifiers,
real-time monitoring, activity logs---becomes a precondition for accountability. Visibility is also
the precondition for noticing drift.

\subsection{From code generation to repository-level agents}

Benchmarks track the shift from advice to action. SWE-bench replaced function-level code generation
with real GitHub issues that must be resolved by editing an entire repository \citep{jimenez2024swebench}.
At its release, the best model solved under 2\% of issues. SWE-agent showed that performance depends
heavily on the \emph{agent--computer interface}, treating the model as a new kind of user whose tools
must be designed for it \citep{yang2024sweagent}. This is an HCI framing applied to the model itself.
OpenHands provided an open platform in which agents write code, use a shell and browse the web in a
sandbox \citep{wang2024openhands}, later re-engineered as an SDK for production agents
\citep{wang2025openhandssdk}. Underneath, most of these systems use the ReAct pattern, which
interleaves reasoning with tool calls \citep{yao2023react}.

Measured in human time, the tasks that frontier models complete with 50\% reliability had reached
around 50 minutes by early 2025, and the trend may have accelerated \citep{kwa2025measuring}. The
authors suggest that the gains are driven primarily by greater reliability and a better ability to
adapt to mistakes, together with better reasoning and tool use. Longer horizons are a
double-edged sword for our question. They mean that agents work for longer stretches without human
input, which is exactly the regime in which an early objective must survive many intervening
observations, tool outputs and sub-tasks.

\subsection{Claude Code as an operational agent}
\label{sec:claudecode}

We use Claude Code \citep{anthropic2026claudecode} as the reference system throughout. It is a
terminal-based agent. In each iteration of its loop the model reads the conversation so far,
proposes one or more tool calls (reading, searching, editing or writing files; running shell
commands; fetching web pages), the harness executes them, and the results are appended to the
context. Six architectural features are directly relevant to goal drift, each with a dual role:
\begin{description}[leftmargin=0pt]
  \item[Task lists.] The agent maintains an explicit to-do list that it updates as it works. A task list
        can anchor the original goal, but it can also \emph{absorb} sub-tasks, giving a distractor the
        same visual and planning status as the goal.
  \item[Sub-agents.] Work can be delegated to sub-agents with their own context window, which
        return only a result. This resembles the planner/coder separation that CodeDelegator uses to
        stop debugging traces from polluting the planner's context \citep{fei2026codedelegator}. At the
        same time, hand-offs are where a drifted state can be inherited \citep{menon2026inherited}.
  \item[Memory files.] Project and user instructions in \texttt{CLAUDE.md} files are loaded into
        every session. Empirically, these manifests contain mostly operational commands,
        implementation notes and architecture \citep{chatlatanagulchai2025manifests}, and context files
        dominate how developers configure agentic tools \citep{galster2026harness}. They can carry a
        persistent ``north star'', but they are also one more channel through which instruction-like
        text enters the context.
  \item[Compaction.] When the context window fills up, the conversation history is summarised.
        Compaction decides what survives of the original request, and the best way to compress is
        still an open question: optimising the compressor's guidelines from observed failures
        improves agents \citep{kang2025acon}, while simply masking old observations matches LLM
        summarisation at lower cost \citep{lindenbauer2025complexity}.
  \item[Hooks.] User-defined commands run deterministically at lifecycle events, for example before
        or after a tool call, when the user submits a prompt, before compaction, or when the agent
        stops. Hooks make it possible to re-inject a goal, to verify requirements before the agent
        declares completion, or to run an external monitor, without relying on the model's
        cooperation.
  \item[Permission modes and headless use.] The user chooses whether each action must be approved,
        whether file edits are auto-accepted, or whether all permissions are bypassed. A non-interactive
        mode emits a machine-readable event stream and allows sessions to be resumed and forked. The
        permission mode is the concrete knob that sets the autonomy level in Feng et al.'s sense. The
        headless mode is what makes controlled experiments possible.
\end{description}

Together, these features turn Claude Code from a conversational assistant into an \emph{operational
entity}: it holds state (files, task list, memory), acts on an external environment, and decides when
it is done. The user's main levers are no longer individual suggestions but the configuration of this
loop: the instructions, the permissions, the hooks and the moments at which they choose to look.
